\documentclass[openany]{ctexbook}
\usepackage{amsmath,amssymb}
\usepackage[hidelinks]{hyperref}
\newcommand{\note}[1]{\footnote{本批译注：#1}}
\title{代数几何基础第五卷预备笔记\\8.16至8.17：校订续译}
\author{据用户上传英译编辑稿译出}
\date{2026年11月15日}
\begin{document}
\maketitle
\chapter*{范围及校订状态}
本批覆盖底本第33页全部可辨文字，包括命题8.16及其证明、备注8.17、随后未编号的残缺论述和术语补记。下一处为第34页命题8.18标题。此为独立增补，未并入原合订本。

底本8.8缺失条件及相关证明缺口仍未恢复。本批审校核对的是翻译、公式及排版，不表示独立验证了所引全部命题。残缺句保留其残缺状态；可确认的记号排印问题订正并说明。
\chapter*{8.16\quad 命题（原稿陈述）}
假设 $f\colon X\to P$ 在一个开子集上非分歧，该开子集的补集余维至少为 $2$；假设 $X$ 几何正规，且在其闭点处的深度至少为 $3$；最后，假设 $X$ 几何整且在 $k$ 上固有。那么，使得 $Y_\xi$ 几何正规、几何整且维数等于 $\dim X-1$ 的点 $\xi\in P^\vee$ 所成集合（是可构造的，并且）其补集的余维至少为 $2$。

\paragraph{原稿证明。}
事实上，由8.14(b)，对于性质“$Y_\xi$ 几何正规且维数为 $\dim X-1$”，上述结论成立（这一维数性质表达的是 $\phi_\xi$ 为 $X_{k(\xi)}$-正则）。\note{8.14前后的假设调整见8.15后的作者补记；其原文缺失和依赖问题仍未全部恢复。此处保留作者引用，不把引用本身作为本批已独立验证的证据。}

再由6.1，所有这些 $Y_\xi$ 都几何连通。由于 $Y_\xi$ 几何正规，它几何整当且仅当它几何连通，由此得到证明。
\section*{8.17\quad 备注}
\begin{enumerate}
\item[(a)] 8.16的假设蕴涵 $\dim X\ge3$。或许只要 $X$ 几何不可约且 $\dim f(X)\ge3$（不要求正规性和非分歧性），使得 $Y_\xi$ 几何不可约的点 $\xi$ 所成集合的补集就具有至少为 $2$ 的余维。\note{底本在此夹有法文“dès que”的转写痕迹。正文译为“只要”，并保留作者“It is possible”的推测语气，不将其升级为定理。}无论如何，我们可以证明这个补集不包含任何超平面（见下文）。\note{原句代词“it”依上下文指例外补集；此处明确所指，不改变其断言范围。}
\item[(b)] 如果去掉闭点 $x$ 处 $\operatorname{prof}_xX\ge3$ 的假设，8.16的结论便不成立。例如，对 $P^3$ 中的一个非奇异二次曲面 $X$，结论就不成立。〔字迹不清〕\footnote{原脚注50：“询问 A.G.”。}这些超平面截面是可约的（事实上由两条相交直线组成），因而形成一个二维族，在 $P^\vee$ 中的余维为 $1$（事实上，该二次曲面的对偶是在对偶空间中由对偶二次型给出的二次曲面……）。\note{底本此处印作“在 $P$ 中”，但所计数的是超平面的参数族，且紧随文字明确提到对偶空间，故订为 $P^\vee$。缺文处未补入“相切”等字样；这里的“这些超平面截面”只承接原稿所讨论的可约截面，不是说该曲面的所有超平面截面都可约。}
\end{enumerate}

对于射影空间中的非奇异曲面，依下一节（或下一条）的讨论，这种情形却应当视为例外情形。假设〔字迹不清〕在 $k$ 上整且固有，并且 $f$ 是浸入。那么由6.1、8.8及8.14可得：如果 $Y^{\mathrm{sing}}\to P^\vee$ 不是泛有限且不可分的，〔原稿接“one of the (???)”，句子残缺〕那些使得 $Y_\xi$ 在 $k(\xi)$ 上可分、并且至多有两个不可约〔字迹不清〕的 $\xi\in P^\vee$，〔句子残缺〕一个余维至少为 $2$ 的补集。\note{本段从假设到结论均有缺字和断句。“not generically finite and inseparable”的否定范围亦不明确；正文保留其不确定性，不自行判定应否定合取，或仅否定“泛有限”。不能由这些残句恢复一个确定的充分条件或结论。8.8亦尚缺完整假设，故不补造推论。“separable”沿用前批译为“可分”。}

现在我们要更仔细地考察曲面的情形（从超平面截面的不可约性这一观点看，曲线的情形显然不成为问题）。\note{“显然”是原作者的表述，底本此处未给证明；本批不另加未经核对的论证。}

（注意：我惊恐地发现，按我一直使用“直纹的”一词时所取的意义，二次曲面竟没有资格称为直纹的。这与前人的用法不一致，有必要为这里使用的概念另造一个更恰当的词。）\note{这是作者对暂用术语的自我修正，不是关于通常意义下二次曲面是否直纹的否定。}
\end{document}
